\section{The CUTLASS reference line}
\label{sec:cutlass}

Comparing a hand written kernel to cuBLAS answers "how close to the vendor", and
that is the headline. It does not answer "how close to a well engineered open
implementation of the same idea", which is the more useful question when the goal
is to find out what my mainloop is missing. So CUTLASS runs as a rung on the same
ladder, through the same harness, on the same shapes~\cite{cutlass}.

It is pinned to release \texttt{v4.7.1}, commit
\texttt{cb4247394dd82148787aed73e5dc7cef33cbf862}, fetched at configure time. The
pin is a tag rather than a branch on purpose: the analysis below is about what
this release's SM120 builders emit, and a moving pin would quietly invalidate it.
Nothing of CUTLASS is vendored into the tree, no CUTLASS target is configured, and
no CUTLASS type appears in a public header. It is headers, included by one
translation unit. The full argument for the 2.x device API over a
\texttt{CollectiveBuilder} on this part is in \texttt{docs/cutlass.md}.

\paragraph{What is established now.} Correctness. The rung is held to the same
three errors as every other tensor rung of Section \ref{sec:accuracy}: $10^{-5}$
relative Frobenius against a same precision cuBLAS oracle, a $\text{tol}(k)$ bound
against a double precision reference over the rounded inputs, and a storage
roundoff bound against the FP32 data those inputs came from. It also passes the
beta equal to zero rule, the empty output and empty contraction cases, and large
K. The sanitizers are clean on it: memcheck, racecheck, initcheck and synccheck
all run without an error on the tensor suite with the rung in it.

\paragraph{What the sweep measured.} The rows exist now, at locked clocks, and
they say the reference line wins everywhere it was run. \texttt{mma\_opt} lands at
91.0 percent of CUTLASS at 1024 cubed, 92.5 at 2048, 93.1 at 4096 and 93.7 at 8192:
a gap that is real, consistent, and narrowing with shape. CUTLASS itself passes
cuBLAS at the two largest shapes, 101.9 percent at 4096 and 101.2 at 8192, which is
worth stating because it means the remaining 6 to 7 points are not a vendor
advantage that a hand written kernel cannot reach. They are engineering I have not
done.

\begin{table}[htbp]
\centering
\begin{tabular}{lrrrr}
\toprule
shape & mma\_opt GFLOP/s & cutlass GFLOP/s & cuBLAS GFLOP/s & mma\_opt as percent of cutlass \\
\midrule
128 cubed & 431 & pending & 745 & pending \\
256 cubed & 2,422 & pending & 4,387 & pending \\
512 cubed & 12,175 & pending & 22,672 & pending \\
768 cubed & 28,368 & pending & 37,648 & pending \\
1024 cubed & 32,025 & 35,209 & 48,210 & 91.0 \\
2048 cubed & 48,559 & 52,506 & 56,370 & 92.5 \\
4096 cubed & 54,602 & 58,669 & 57,568 & 93.1 \\
8192 cubed & 56,375 & 60,154 & 59,450 & 93.7 \\
\bottomrule
\end{tabular}

\caption{The gap against the CUTLASS reference line, from the committed locked
clock sweep. CUTLASS was not run at 128 to 768 cubed, and those cells read pending
rather than being extrapolated from the shapes that were run.}
\label{tab:cutlassgap}
\end{table}

\paragraph{Where the gap is, now that cuBLAS has been profiled.} Round 13 settled
part of this by profiling the vendor directly, and it corrected the prediction this
chapter was written against. The expectation was Ampere era
\texttt{cutlass\_80\_tensorop\_h16816gemm} kernels. What cuBLAS actually dispatches
on this part is an \texttt{nvjet} kernel compiled for this architecture, and four
things in it are the gap. It stages \emph{both} operands through TMA where
\texttt{mma\_opt} uses \texttt{cp.async}. It picks non power of two N tiles, 176 at
1024 and 80 at 4096, which the six power of two shapes in this tree cannot express.
It spends 255 registers per thread against 122, running at 16.67 percent
theoretical occupancy against 33.33, so it does not try to hide latency with warps
at all. And it takes up to 99.33 KB of dynamic shared memory per block, the full
\texttt{sharedMemPerBlockOptin} budget, against 49.15 KB.

One candidate is now eliminated rather than open: bank conflicts are not where the
gap lives. The vendor kernels measure 33,792 conflicts at 1024, 527,360 at 4096 and
2,097,152 at 8192, against 106,296 and 465,929 for \texttt{mma\_opt} at 4096 and
8192. The hand kernel has the cleaner shared path of the two by that counter and is
still slower, so the shared memory layout is not the thing to spend the next round
on.

The remaining candidates stay hypotheses. CUTLASS runs 128 threads with 64 by 64
warp tiles while \texttt{mma\_opt} runs 256 threads with 64 by 32, so the wider
warp tile amortizes each \texttt{ldmatrix} over more \texttt{mma.sync} at the cost
of accumulator registers; the HMMA to LDSM ratio in the SASS captures is already
on file to test that. CUTLASS stages its accumulator through shared memory for
coalesced global stores while \texttt{mma\_opt} stores 64 bit pairs straight from
the fragment layout. Both run three stages, but the barrier placement and the
depth of the register prefetch differ. And the reference line takes 236 registers
over 128 threads against \texttt{mma\_opt}'s 122 over 256, with neither spilling,
so which of the two the machine actually lets sit resident is the question Section
\ref{sec:sass} answers for one of them and not yet for the other.
